\begin{figure}[htbp]
\centering
\setlength{\unitlength}{1pt}
\begin{picture}(370,128)
% Opposite shaded quadrants are the dielectric sectors in the sharp model.
\put(24,29){\color{black!20}\rule{35pt}{35pt}}
\put(59,64){\color{black!20}\rule{35pt}{35pt}}
\put(24,29){\framebox(70,70){}}
\put(59,29){\line(0,1){70}}
\put(24,64){\line(1,0){70}}
\put(59,64){\circle{54}}
\put(59,64){\circle*{3}}
\put(11,107){\makebox(100,13){sharp local model}}
\put(27,6){\makebox(64,13){one corner}}
\put(106,77){\vector(-1,-1){34}}
\put(108,77){\makebox(80,12)[l]{observation disk}}
\put(138,64){\vector(-1,0){76}}
\put(141,59){\makebox(100,14)[l]{meeting point}}
\put(250,90){\makebox(110,13)[l]{shaded: dielectric}}
\put(250,69){\makebox(110,13)[l]{white: background}}
\put(250,43){\makebox(110,13)[l]{circle: disk boundary}}
\end{picture}
\caption{Schematic of one sharp two-dimensional checkerboard neighborhood.
Two opposite dielectric sectors meet at a point; the other sectors are
background. The circle indicates a fixed observation disk, not a material
interface. In the proof finitely many disjoint copies are illuminated at
once. The final interfaces are rounded and smooth; the sharp meeting point
is used only as a reference for constructing their peaks.}
\label{hs:fig-corner}
\end{figure}
